\section{Blocking sets and projective codes}\label{sec:translations}

The complement of a \((k,n)\)-arc \(\mathcal K\) in a projective plane of
order \(q\) is a \((q+1-n)\)-fold blocking set.  Completeness says more: each
point of the complement lies on a line meeting it in exactly \(q+1-n\)
points.  Thus the complement is a minimal multiple blocking set.  Under this
correspondence, \(n\)-secants of \(\mathcal K\) are the
\((q+1-n)\)-secants of its complement, and the hypothesis
\(\sigma_n(Q)\ge\lambda\) says that every point of the blocking set is on at
least \(\lambda\) such secants.  This is precisely the dual setting of
Bishnoi--Mattheus--Schillewaert~\cite[Section~8]{BMS}.

In particular, Theorem~\ref{thm:characteristic-three} is equivalent to the
following statement.

\begin{corollary}\label{cor:blocking-set-form}
\coverage{absent}
Let \(q=3^h\) and let \(B\) be the complement of a complete
\((k,2q/3+1)\)-arc in \(\PG(2,q)\).  Then \(B\) is a minimal
\(q/3\)-fold blocking set and
\[
 |B|\le \frac{2q^2}{3}-\frac q3+1+o(q).
\]
\uses{thm:characteristic-three}
\end{corollary}

There is also a standard coding interpretation; compare Kohnert's incidence
criterion for projective code extension~\cite{Kohnert}.  Projective representatives
of the points of a \((k,n)\)-arc form the columns of a generator matrix for a
projective \([k,3,k-n]_q\) code.  A line meeting the arc in \(n\) points
corresponds to a minimum-weight projective codeword.  Adding a new projective
column raises the minimum distance by one exactly when no minimum-weight
hyperplane contains that column.  Hence completeness is equivalent to
nonextendibility.  The condition \(\sigma_n(Q)\ge\lambda\) records at least
\(\lambda\) minimum-weight hyperplanes through every candidate column.

The same statement is available in every projective dimension.  Points and
hyperplanes of \(\PG(r-1,q)\) form a symmetric design with parameters
\begin{equation}\label{eq:projective-design-parameters}
 v=\frac{q^r-1}{q-1},\qquad
 K=\frac{q^{r-1}-1}{q-1},\qquad
 \Lambda=\frac{q^{r-2}-1}{q-1}.
\end{equation}

\begin{corollary}[Projective codes]\label{cor:projective-codes}
\coverage{absent}
Let \(C\) be a projective \([k,r,k-n]_q\) code, and let \(\cS\) be its
minimum-weight hyperplanes.  Partition the points of \(\PG(r-1,q)\) into the
\(k\) columns of \(C\) and the remaining candidate columns.  For every
integer \(T=|\cS|\), the two minimum-hyperplane degree sequences satisfy the
exact interval condition \eqref{eq:interval-overlap}, with the parameters in
\eqref{eq:projective-design-parameters}.  Any prescribed lower bound on the
number of minimum-weight hyperplanes through a candidate column may be used
as the lower endpoint of its integer degree interval.
\end{corollary}

\begin{proof}
Apply Theorem~\ref{thm:paired-envelope} to the point--hyperplane design of
\(\PG(r-1,q)\).  A minimum-weight hyperplane contains exactly \(n\) columns,
so the selected blocks have constant intersection size \(n\) with the first
point class.
\uses{thm:paired-envelope}
\end{proof}

Exact two-valued hyperplane intersections correspond to projective two-weight
codes~\cite{CalderbankKantor}.  In the simultaneously tight families of
Proposition~\ref{prop:simultaneous-equality}, all but \(O(q)\) dual lines have
the two intersection numbers
\[
 uv\quad\hbox{and}\quad (u+1)(v+1).
\]
Their difference is \(u+v+1\), which must be a power of the characteristic
along an infinite Desarguesian family.  This agrees with the classical
divisibility condition for exact projective two-weight codes.  The present
argument does not classify the \(O(q)\) exceptional lines.

\section{Concluding remarks}\label{sec:conclusion}

The first two incidence moments already contain the classical spectral bound.
Their sharp integer degree ranges retain information that the real relaxation
loses, and at each rational equality parameter this loss has linear order.
Modular stability then converts the nearly two-valued dual intersection
distribution into a bounded repair.  Completeness assigns a further linear
cost to every point in that repair support.

The characteristic-three result leaves a concrete extremal problem:
\begin{equation}\label{eq:open-asymptotic}
 t_{2q/3+1}(2,q)
 \stackrel{?}{=}\frac{q^2}{3}+\frac{4q}{3}+o(q)
 \qquad(q=3^h\to\infty).
\end{equation}
A matching construction would determine the first two asymptotic terms.  A
second problem is structural: classify the sets of dual points with only
\(O(q)\) lines outside the two predicted intersection numbers and determine
which of them can occur as the full set of \(n\)-secants of a complete
\((k,n)\)-arc.
